% Options for packages loaded elsewhere
\PassOptionsToPackage{unicode}{hyperref}
\PassOptionsToPackage{hyphens}{url}
\PassOptionsToPackage{dvipsnames,svgnames,x11names}{xcolor}
%
\documentclass[
  11pt,
]{article}
\usepackage{amsmath,amssymb}
\usepackage{iftex}
\ifPDFTeX
  \usepackage[T1]{fontenc}
  \usepackage[utf8]{inputenc}
  \usepackage{textcomp} % provide euro and other symbols
\else % if luatex or xetex
  \usepackage{unicode-math} % this also loads fontspec
  \defaultfontfeatures{Scale=MatchLowercase}
  \defaultfontfeatures[\rmfamily]{Ligatures=TeX,Scale=1}
\fi
\usepackage{lmodern}
\ifPDFTeX\else
  % xetex/luatex font selection
\fi
% Use upquote if available, for straight quotes in verbatim environments
\IfFileExists{upquote.sty}{\usepackage{upquote}}{}
\IfFileExists{microtype.sty}{% use microtype if available
  \usepackage[]{microtype}
  \UseMicrotypeSet[protrusion]{basicmath} % disable protrusion for tt fonts
}{}
\makeatletter
\@ifundefined{KOMAClassName}{% if non-KOMA class
  \IfFileExists{parskip.sty}{%
    \usepackage{parskip}
  }{% else
    \setlength{\parindent}{0pt}
    \setlength{\parskip}{6pt plus 2pt minus 1pt}}
}{% if KOMA class
  \KOMAoptions{parskip=half}}
\makeatother
\usepackage{xcolor}
\usepackage[margin=1in]{geometry}
\usepackage{longtable,booktabs,array}
\usepackage{calc} % for calculating minipage widths
% Correct order of tables after \paragraph or \subparagraph
\usepackage{etoolbox}
\makeatletter
\patchcmd\longtable{\par}{\if@noskipsec\mbox{}\fi\par}{}{}
\makeatother
% Allow footnotes in longtable head/foot
\IfFileExists{footnotehyper.sty}{\usepackage{footnotehyper}}{\usepackage{footnote}}
\makesavenoteenv{longtable}
\setlength{\emergencystretch}{3em} % prevent overfull lines
\providecommand{\tightlist}{%
  \setlength{\itemsep}{0pt}\setlength{\parskip}{0pt}}
\setcounter{secnumdepth}{-\maxdimen} % remove section numbering
\ifLuaTeX
  \usepackage{selnolig}  % disable illegal ligatures
\fi
\IfFileExists{bookmark.sty}{\usepackage{bookmark}}{\usepackage{hyperref}}
\IfFileExists{xurl.sty}{\usepackage{xurl}}{} % add URL line breaks if available
\urlstyle{same}
\hypersetup{
  pdfauthor={Stefano Alimonti --- stefalimonti@gmail.com},
  colorlinks=true,
  linkcolor={blue},
  filecolor={Maroon},
  citecolor={Blue},
  urlcolor={blue},
  pdfcreator={LaTeX via pandoc}}

\author{Stefano Alimonti --- stefalimonti@gmail.com}
\date{March 2026}

\begin{document}

{
\hypersetup{linkcolor=black}
\setcounter{tocdepth}{2}
\tableofcontents
}
\section{π from Pure Arithmetic: A Spectral Phase Transition in the Binary Carry
Bridge}\label{ux3c0-from-pure-arithmetic-a-spectral-phase-transition-in-the-binary-carry-bridge}

\textbf{Author:} Stefano Alimonti \textbf{Affiliation:} Independent Researcher
\textbf{Date:} March 2026

\begin{center}\rule{0.5\linewidth}{0.5pt}\end{center}

\newpage

\subsection{Abstract}\label{abstract}

We study the propagation of arithmetic carries in binary multiplication and
discover a spectral phase transition that converts rational arithmetic into a
transcendental constant. Given two uniformly random K-bit integers X, Y with
product P = XY written in D = 2K − 1 bits, define the \emph{sector ratio} R(K) =
σ₁₀(K)/σ₀₀(K), where σ\_ac is the carry-weighted partition function restricted
to pairs whose second-highest bits of X, Y are (a, c).

We prove:

\begin{enumerate}
\def\labelenumi{\arabic{enumi}.}
\tightlist
\item
  The Markov (independent-convolution) model gives R\_Markov = +2/3, a rational
  constant independent of K (Theorem 2).
\item
  The alternating spectral sum over the Dirichlet bridge has the exact closed
  form S(β) = 2β/(1 + 2β), a Möbius transformation (Theorem 1).
\item
  Setting S = −π yields the critical parameter A* = (3π + 2)/(2(π + 1)), placing
  the system at distance 1/(2(π + 1)) from the pole (Corollary 2).
\item
  The spectral zeta function of the Markov carry bridge reproduces ζ(2s) at all
  even arguments of the Riemann zeta function, via Weyl asymptotics (Theorem 3).
\item
  No D-odd pair has sector (1,1) (Theorem 4), and the count ratio N₁₀/N₀₀ →
  (2ln(4/3) − 1/2)/(2ln(9/8)) = 0.31993\ldots{} involves only ln 2 and ln 3
  (Theorem 5). The entire factor π in R enters exclusively through the carry
  weight ratio ⟨val⟩₁₀/⟨val⟩₀₀.
\end{enumerate}

We conjecture:

\begin{enumerate}
\def\labelenumi{\arabic{enumi}.}
\setcounter{enumi}{5}
\tightlist
\item
  The exact cascade sector ratio satisfies R(∞) = lim\_\{K→∞\} R(K) = −π = −4
  L(1, χ₄) (Conjecture 1). Exact enumeration up to K = 21 (1.1 × 10¹² digit
  pairs) confirms this to 4.0 significant digits by exponential Richardson
  extrapolation (§9.2; 4.4 digits using exact values only). The weight sin(nπ/2)
  in the spectral sum equals the Dirichlet character χ₄(n), which emerges as the
  effective spectral weight under resolvent universality {[}E{]}.
\end{enumerate}

The system undergoes a spectral phase transition: the Markov model (subcritical)
produces the rational +2/3; digit correlations push the carry bridge past a
critical point into a supercritical regime where the transcendental −π emerges.
Combined with ζ(2s) from the spectral zeta function and L(1, χ₄) from the sector
ratio, the carry bridge produces both factors of the Dedekind zeta function
ζ\_\{ℚ(i)\}(s) = ζ(s) · L(s, χ₄) of the Gaussian integers.

\begin{center}\rule{0.5\linewidth}{0.5pt}\end{center}

\newpage

\subsection{1. Introduction}\label{introduction}

The appearance of π in number theory is usually traced to analytic methods: the
residue calculus in the proof of the prime number theorem, the Fourier analysis
underlying Dirichlet L-functions, or the Gaussian integral in the central limit
theorem. In each case, the geometric content of π (circumference, area,
rotation) enters through an analytical tool applied to a discrete problem.

In this paper we exhibit a mechanism that produces π from \emph{purely
arithmetic operations} --- binary multiplication and carry propagation ---
without invoking any geometric or analytic machinery except at the final step of
evaluating a closed-form identity. The mechanism reveals a spectral phase
transition: a rational quantity (+2/3) is converted into a transcendental one
(−π) when digit correlations push the system past a critical point.

We also show that the Markov transfer operator for the carry bridge, viewed as a
spectral zeta function, reproduces ζ(2s) at all even arguments --- a
manifestation of Weyl's law in the arithmetic setting. Combined with the
conjectural L(1, χ₄) = π/4 from the sector ratio, this connects the carry bridge
to the Dedekind zeta function of the Gaussian integers.

\textbf{Scope and status.} Five theorems are proved unconditionally (Theorems
1--5, Theorem 6, Theorem 7). The central claim R(∞) = −π (Conjecture 1) is
supported by 4.0-digit numerical evidence (exponential Richardson, §9.2) but
remains open; under the Linear Mix Hypothesis of {[}E{]}, it follows from the
proved closed-form identity S(A) = 2(1−A)/(3−2A). This paper is therefore a
contribution to experimental mathematics: it identifies the phenomenon, proves
the structural framework, and isolates the single remaining spectral conjecture.
The remaining proof targets are: (i) prove the scalar identity C = −4 in R(∞) =
C · L(1, χ₄), and (ii) prove analytically that the conditioned D-odd chain
preserves the dominant 1/2 spectral rate. The stopping-time framework extends to
a full Dirichlet series in the complex variable s in {[}L{]}, revealing L²(s,
χ₄) structure.

\subsubsection{1.1 Setting}\label{setting}

Let X and Y be independent uniformly random integers in {[}2\^{}\{K−1\}, 2\^{}K)
(K-bit integers with leading bit 1). Their product P = XY has D = 2K − 1 binary
digits. The schoolbook multiplication algorithm computes

\begin{verbatim}
P = Σ_{j=0}^{D-1} p_j · 2^j
\end{verbatim}

where at each position j the product digit and carry are determined by

\begin{verbatim}
p_j = (conv_j + c_j) mod 2,    c_{j+1} = ⌊(conv_j + c_j) / 2⌋
\end{verbatim}

with conv\_j = Σ\_i x\_i · y\_\{j−i\} the convolution at position j and c\_0 = 0
the boundary condition.

\textbf{Definition 1 (Sector partition).} The \emph{sector} of a pair (X, Y) is
(a, c) ∈ \{0,1\}² where a = x\_\{K−2\} and c = y\_\{K−2\} are the second-highest
bits. The \emph{sector partition function} is

\begin{verbatim}
σ_ac(K) = Σ_{(X,Y): x_{K-2}=a, y_{K-2}=c} w(X, Y)
\end{verbatim}

where \(w(X,Y) = w_{\mathrm{casc}}(X,Y)\) is the cascade carry weight defined
below. Two valuations appear in the literature: the \emph{schoolbook} weight
\(w_{\mathrm{sch}} = 2c_{D-2} - 1\) (fixed-position readout) and the
\emph{cascade} weight \(w_{\mathrm{casc}} = c_{M-1} - 1\) where \(M\) is the
topmost nonzero carry position (first-passage readout; not to be confused with
the companion matrix \(M\) of {[}A{]}, {[}B{]}). Throughout this paper
\(w = w_{\mathrm{casc}}\) unless otherwise noted; it is this valuation that
produces \(-\pi\).

\textbf{Definition 2 (Sector ratio).} R(K) = σ₁₀(K)/σ₀₀(K). By symmetry, σ₁₀ =
σ₀₁. (Not to be confused with the carry correction factor \(R(l,s)\) of {[}B{]},
which measures the per-prime deviation from the Euler product.)

\subsubsection{1.2 Main results}\label{main-results}

\textbf{Theorem 1 (Spectral closed form).} For β ∈ (−1/2, +∞), define

\begin{verbatim}
S(β, L) = Σ_{k=0}^{⌊L/2⌋-1} (-1)^k / (1/2 + β sin²((2k+1)π/(2L)))
\end{verbatim}

Then

\begin{verbatim}
lim_{L→∞} S(β, L) = 2β/(1 + 2β)
\end{verbatim}

\textbf{Corollary 1 (Möbius form).} Setting β = 1 − A, the sum as a function of
the eigenvalue shift A is

\begin{verbatim}
S(A) = 2(1−A)/(3−2A)
\end{verbatim}

a Möbius transformation with a zero at A = 1 and a pole at A = 3/2.

\textbf{Corollary 2 (Phase transition).}

\begin{itemize}
\tightlist
\item
  A = 0 (Markov): S = +2/3
\item
  A = 1 (critical): S = 0
\item
  A = A\emph{: S = −π, where A} = (3π + 2)/(2(π + 1)) = 3/2 − 1/(2(π+1)) ≈
  1.37927
\end{itemize}

\textbf{Theorem 2 (Markov baseline).} The independent-convolution (Markov) model
gives R\_Markov = +2/3, with finite-size correction

\begin{verbatim}
S(0, L) = 2/3 + 10 ζ(2)/(3 L²) + O(L⁻⁴)
\end{verbatim}

\textbf{Theorem 3 (Spectral zeta function).} Let F = ½I − T\_bridge be the
fluctuation operator of the Markov carry bridge of length L, with eigenvalues
ε\_n = sin²(nπ/(2(L+1))). Then for Re(s) \textgreater{} 1/2,

\begin{verbatim}
ζ_F(s) / (2L/π)^{2s}  →  ζ(2s)    as L → ∞
\end{verbatim}

with convergence rate \(O(L^{1-2\sigma})\) where \(\sigma = \Re(s)\) (reducing
to \(O(1/L^2)\) for \(\sigma > 3/2\)).

\textbf{Theorem 4 (N₁₁ = 0).} No D-odd pair has sector (1,1).

\textbf{Theorem 5 (Closed-form count ratio).} N₁₀/N₀₀ → (2ln(4/3) −
1/2)/(2ln(9/8)) = 0.31993\ldots, involving only ln 2 and ln 3. The convergence
correction is O(1/2\^{}K).

\textbf{Conjecture 1 (Exact sector ratio).} R(∞) = lim\_\{K→∞\} σ₁₀(K)/σ₀₀(K) =
−π. (Note: this is implied by, but weaker than, the Linear Mix Hypothesis of
{[}E{]}, which is also called ``Conjecture 1'' in that paper.)

Supported by exact enumeration to K = 21 with Richardson extrapolation (Section
9). By the decomposition R = R₀ + ΔR {[}E{]}, the conjecture reduces to proving
that the cascade correction ΔR = Tr{[}(I−T)⁻¹·P\_\{c=0\}{]} equals exactly −π −
R₀ = +0.7896\ldots{}

\subsubsection{1.3 Structure of the paper}\label{structure-of-the-paper}

\begin{itemize}
\tightlist
\item
  Section 2: Carry bridge and Dirichlet eigenfunctions
\item
  Section 3: Markov baseline (R = +2/3)
\item
  Section 4: Proof of the closed form (Theorem 1)
\item
  Section 5: The spectral phase transition
\item
  Section 6: The Dirichlet character χ₄ and L(1, χ₄)
\item
  Section 7: The spectral zeta function and ζ(2s) (Theorem 3)
\item
  Section 8: The Bernoulli hierarchy and the Gaussian integers
\item
  Section 9: Numerical evidence
\item
  Section 10: The gap between S(A) and R(K)
\item
  Section 11: Open problems
\end{itemize}

\begin{center}\rule{0.5\linewidth}{0.5pt}\end{center}

\newpage

\subsection{2. The Carry Bridge and Dirichlet
Eigenfunctions}\label{the-carry-bridge-and-dirichlet-eigenfunctions}

\subsubsection{2.1 The D-odd constraint}\label{the-d-odd-constraint}

A pair (X, Y) with X, Y ∈ {[}2\^{}\{K−1\}, 2\^{}K) has product P with exactly D
= 2K − 1 bits (the ``D-odd'' constraint). This forces c\_0 = 0 (no carry into
position 0) and c\_\{D−1\} = 0 (no overflow beyond D − 1 bits). The carry chain
c\_0, c\_1, \ldots, c\_\{D−1\} with these boundary conditions forms a
\emph{carry bridge} --- a lattice path conditioned to return to zero, analogous
to a Brownian bridge.

\subsubsection{2.2 Transfer operator and
eigenvalues}\label{transfer-operator-and-eigenvalues}

The Markov transfer operator for the carry chain, conditioned on independent
convolutions at each position, was introduced by Diaconis and Fulman {[}1{]}. In
the base-2 case, the carry at each position takes values in \{0, 1\}, and the
transition matrix has eigenvalues

\begin{verbatim}
λ_n = (1/2) cos(nπ/L),    n = 1, …, L−1
\end{verbatim}

where L = D = 2K − 1 is the bridge length. The corresponding eigenfunctions are
the Dirichlet sine functions φ\_n(j) = sin(nπj/L), reflecting the zero boundary
conditions c\_0 = c\_L = 0. The leading eigenvalue λ\_1 = (1/2)cos(π/L) → 1/2
encodes the Diaconis--Fulman spectral gap.

\subsubsection{2.3 Spectral representation of the sector
response}\label{spectral-representation-of-the-sector-response}

The sector perturbation at position j* = K − 2 excites the Dirichlet modes of
the carry bridge. The response at the boundary is a spectral sum

\begin{verbatim}
R ~ Σ_{n=1}^{L-1} sin(nπj*/L) / (1 − λ_n^eff)
\end{verbatim}

where λ\_n\^{}eff incorporates both the Markov eigenvalue and the effective
modification due to digit correlations. In the Markov model, the perturbation
position satisfies j\emph{/L = (K−2)/(2K−1) → 1/2 as K → ∞, giving sin(nπj}/L) →
sin(nπ/2) = χ₄(n), the non-principal Dirichlet character modulo 4 (§6). In the
true cascade, the stopping position is not localized at the midpoint (it
concentrates near the top of the chain; see {[}E, §8.5{]}). Nonetheless, under
the LMH, the resolvent-weighted spectral sum converges to −π through a
collective mechanism {[}E, §8.3{]} that is functionally equivalent to the
midpoint evaluation (§6.2).

The precise relationship between R(K) and the spectral sum S(A, L) is subtle:
{[}E, §8.1--§8.3{]} establish that S(A, L) is not a direct model of R(K), but
rather that both quantities converge (numerically) toward −π through different
mechanisms (§10). Exact enumeration through K = 21 (\(5.5 \times 10^{11}\)
pairs; {[}E, Tables 1--2{]}) confirms convergence to −π with 2.1 significant
digits raw and 4.0 digits by exponential Richardson extrapolation (§9). {[}E{]}
identifies the mechanism as \textbf{resolvent universality} of the effective
transfer operator. Formalizing this as a rigorous identity remains the principal
open problem.

\begin{center}\rule{0.5\linewidth}{0.5pt}\end{center}

\newpage

\subsection{3. The Markov Baseline: R = +2/3}\label{the-markov-baseline-r-23}

\subsubsection{3.1 Independent convolutions}\label{independent-convolutions}

In the Markov model, convolutions at distinct positions are treated as
independent random variables given the carry state. The sector perturbation
modifies the convolution distribution at position K − 2 by conditioning on the
second-highest bits (a, c) of the input.

\subsubsection{3.2 Proof of S(0, L) → 2/3}\label{proof-of-s0-l-23}

With A = 0 (no correlation-induced shift), the spectral sum reduces to

\begin{verbatim}
S(0, L) = Σ_{n odd} (-1)^{(n-1)/2} / (1 − (1/2)cos(nπ/L))
\end{verbatim}

This is S(β = 1, L) in the notation of Theorem 1, giving S(β = 1) = 2·1/(1+2) =
2/3.

\textbf{Proposition 1 (Finite-size correction).}

\begin{verbatim}
S(0, L) − 2/3 = 5π²/(9L²) + O(L⁻⁴) = 10 ζ(2)/(3L²) + O(L⁻⁴)
\end{verbatim}

The correction coefficient 5π²/9 is derived analytically from the Fourier
expansion (Section 4) and verified numerically to 6 significant digits.

\begin{center}\rule{0.5\linewidth}{0.5pt}\end{center}

\newpage

\subsection{4. Proof of the Closed Form}\label{proof-of-the-closed-form}

\textbf{Proof of Theorem 1.} The proof proceeds in five steps.

\textbf{Step 1: Fourier expansion.} The integrand has the standard Fourier
series (see e.g.~{[}10{]})

\begin{verbatim}
1/(1/2 + β sin²t) = (2/γ) [1 + 2 Σ_{m=1}^∞ r^m cos(2mt)]
\end{verbatim}

where γ = √(1 + 2β) and r = (γ − 1)²/(2β), valid for \textbar r\textbar{}
\textless{} 1 (equivalently β \textgreater{} −1/2). This follows from the
identity for 1/(A − B cos θ) with A = (1 + β), B = β, θ = 2t.

\textbf{Step 2: Alternating midpoint sum.} For L ≡ 0 (mod 4), the alternating
midpoint sum satisfies

\begin{verbatim}
T_m(L) := Σ_{k=0}^{L/2-1} (-1)^k cos(m(2k+1)π/L) =
    1/cos(mπ/L)   if m is odd,
    0              if m is even.
\end{verbatim}

\emph{Proof.} By geometric summation, Σ\_k (-1)\^{}k e\^{}\{im(2k+1)π/L\} =
e\^{}\{imπ/L\}(1 − (−1)\^{}\{L/2+m\}) / (1 + e\^{}\{2imπ/L\}). When L/2 is even
and m is odd, the numerator is 2 and the denominator is 2cos(mπ/L). ∎

\textbf{Step 3: Assembly.} Substituting the Fourier expansion into S(β, L):

\begin{verbatim}
S(β, L) = (2/γ) [T_0 + 2 Σ_{m=1}^∞ r^m T_m(L)]
        = (4/γ) Σ_{j=0}^∞ r^{2j+1} / cos((2j+1)π/L)
\end{verbatim}

(T\_0 = 0 because the alternating sum of a constant vanishes when L/2 is even.)

\textbf{Step 4: Continuum limit.} As L → ∞, cos((2j+1)π/L) → 1 for each fixed j.
Since \textbar r\textbar{} \textless{} 1, dominated convergence gives

\begin{verbatim}
S(β, ∞) = (4r/γ) · 1/(1 − r²)
\end{verbatim}

\textbf{Step 5: Simplification.} Direct computation using r = (γ − 1)²/(2β) and
γ² = 1 + 2β:

\begin{verbatim}
4r / (γ(1 − r²)) = 4·(γ−1)²/(2β) / (γ · 4γ(γ−1)/(2β)²)
\end{verbatim}

After cancellation, this simplifies to

\begin{verbatim}
S(β, ∞) = 2β/(1 + 2β)   ∎
\end{verbatim}

\textbf{Proposition 2 (General correction coefficient).} The O(1/L²) correction
to S(β, L) is

\begin{verbatim}
S(β, L) = 2β/(1 + 2β) + C(β)/L² + O(L⁻⁴)
\end{verbatim}

where

\begin{verbatim}
C(β) = 2π² r(1 + 6r² + r⁴) / (γ(1 − r²)³)
\end{verbatim}

For β = 1: C(1) = 5π²/9 = 10ζ(2)/3, confirming Proposition 1.

\emph{Proof.} Expand 1/cos((2j+1)π/L) = 1 + (2j+1)²π²/(2L²) + O(L⁻⁴), substitute
into Step 3, and sum the geometric series. ∎

\textbf{Remark (M-parity subtlety).} The closed form S(β,∞) = 2β/(1+2β) applies
when M = ⌊L/2⌋ is even. For M odd, the limit differs by an O(1) term at finite
L, but both subsequences converge to the same limit. This is because the
additional term vanishes as L → ∞.

\begin{center}\rule{0.5\linewidth}{0.5pt}\end{center}

\newpage

\subsection{5. The Spectral Phase
Transition}\label{the-spectral-phase-transition}

\subsubsection{5.1 Interpretation}\label{interpretation}

The closed form S(A) = 2(1−A)/(3−2A) reveals a \emph{spectral phase transition}
parameterized by the eigenvalue shift A:

\begin{longtable}[]{@{}lll@{}}
\toprule\noalign{}
A & S(A) & Regime \\
\midrule\noalign{}
\endhead
\bottomrule\noalign{}
\endlastfoot
0 & +2/3 & Subcritical (Markov): rational \\
1/2 & +1/2 & Subcritical \\
1 & 0 & Critical: exact cancellation \\
1.2 & −2/3 & Supercritical \\
A* ≈ 1.379 & −π & Supercritical: transcendental \\
3/2 & ±∞ & Pole (divergence) \\
\end{longtable}

\subsubsection{5.2 The critical point}\label{the-critical-point}

At A = 1 (β = 0), all effective eigenvalues equal 1/2, and every term in the
alternating sum has the same magnitude. The sum vanishes by exact cancellation:
Σ (−1)\^{}k · 2 = 0. This is the boundary between the subcritical (positive) and
supercritical (negative) regimes.

\subsubsection{5.3 Subcritical
vs.~supercritical}\label{subcritical-vs.-supercritical}

For A \textless{} 1, the effective eigenvalues are below 1/2 and the alternating
sum is positive (dominated by the k = 0 term). For A \textgreater{} 1,
eigenvalues exceed 1/2, denominators shrink, and the negative terms dominate.
The sign flip at A = 1 is the \emph{trace anomaly}: digit correlations reverse
the sign of the classical (Markov) result.

\subsubsection{5.4 Distance to the pole}\label{distance-to-the-pole}

The value A* = 3/2 − 1/(2(π+1)) places the physical system at distance
1/(2(π+1)) ≈ 0.121 from the pole of the Möbius transformation. The formula is
self-referential: A* contains π, and S(A*) = −π. This is not circular --- it is
a well-defined algebraic relation 2(1−A)/(3−2A) = −π with a unique solution A =
(3π+2)/(2(π+1)).

\begin{center}\rule{0.5\linewidth}{0.5pt}\end{center}

\newpage

\subsection{6. The Dirichlet Character and
L-function}\label{the-dirichlet-character-and-l-function}

\subsubsection{6.1 sin(nπ/2) = χ₄(n)}\label{sinnux3c02-ux3c7ux2084n}

The weight function in the spectral sum takes the values

\begin{verbatim}
sin(nπ/2) = +1 if n ≡ 1 (mod 4),
           = −1 if n ≡ 3 (mod 4),
           =  0 if n is even,
\end{verbatim}

which is exactly the non-principal Dirichlet character χ₄: (ℤ/4ℤ)* → \{±1\}.

\subsubsection{6.2 The effective χ₄
model}\label{the-effective-ux3c7ux2084-model}

In the Markov (independent-convolution) model, the sector perturbation acts at
position j* = K − 2, and the bridge has length L = 2K − 1. Since j\emph{/L → 1/2
as K → ∞, the Dirichlet mode evaluation sin(nπj}/L) → sin(nπ/2) = χ₄(n). This
``midpoint selection'' provides a clean motivation for why χ₄ appears in the
spectral sum.

\emph{Caveat ({[}E{]}).} In the true (non-Markov) cascade, the stopping position
concentrates near the top of the chain (depth j = 2, 3, 4 from position D−2),
not at the midpoint. The physical cascade stopping profile, projected onto the
Dirichlet sine basis, does not equal χ₄(n) at finite K {[}E, §8{]}. Nonetheless,
the resolvent

\[(I - \mathcal{K}_{\mathrm{eff}})^{-1}\]

acts as a matched spectral filter that transmutes the true erratic cascade
profiles into a macroscopic sum converging (under LMH) to −π {[}E, §8.3{]}. The
χ₄ model is therefore best understood as an \emph{effective toy reward} --- a
mathematically tractable proxy that belongs to the same universality class as
the true cascade in the thermodynamic limit.

\subsubsection{6.3 Connection to L(1, χ₄)}\label{connection-to-l1-ux3c7ux2084}

The Dirichlet L-function at s = 1 gives the Leibniz formula:

\begin{verbatim}
L(1, χ₄) = Σ_{n=0}^∞ (−1)^n/(2n+1) = 1 − 1/3 + 1/5 − ⋯ = π/4
\end{verbatim}

If R = −π (Conjecture 1), then R = −4 L(1, χ₄). The carry bridge produces π
through the same Dirichlet series that appears in the classical evaluation of π
via the Leibniz--Gregory formula.

\subsubsection{6.4 Structural identity}\label{structural-identity}

Combined with the sub-sector ratio α → 1/6 = B₂ (the second Bernoulli number)
and ζ(2) = π²/6:

\begin{verbatim}
α · R² = (1/6) · π² = π²/6 = ζ(2)
\end{verbatim}

If both α = 1/6 (Conjecture 2, §8) and R = −π (Conjecture 1) can be established
from carry arithmetic, this would yield a new proof of the Basel problem --- one
that begins with integer multiplication. \textbf{Caveat:} this depends on two
currently unproved conjectures.

\begin{center}\rule{0.5\linewidth}{0.5pt}\end{center}

\newpage

\subsection{7. The Spectral Zeta Function and
ζ(2s)}\label{the-spectral-zeta-function-and-ux3b62s}

\subsubsection{7.1 The fluctuation operator}\label{the-fluctuation-operator}

\textbf{Definition 3 (Fluctuation operator).} F = ½I − T\_bridge, where
T\_bridge is the Markov carry bridge transfer matrix. The eigenvalues of F are

\begin{verbatim}
ε_n = 1/2 − λ_n = sin²(nπ/(2(L+1))),    n = 1, …, L−1
\end{verbatim}

\textbf{Definition 4 (Spectral zeta function).}

\begin{verbatim}
ζ_F(s) = Σ_{n=1}^{L-1} ε_n^{-s} = Σ_{n=1}^{L-1} sin^{-2s}(nπ/(2(L+1)))
\end{verbatim}

\subsubsection{7.2 Proof of Theorem 3}\label{proof-of-theorem-3}

For n ≪ L, the small-angle approximation sin(x) ≈ x gives ε\_n ≈ (nπ/(2L))², so

\begin{verbatim}
ζ_F(s) ≈ (2L/π)^{2s} Σ_{n=1}^{L-1} n^{-2s}
\end{verbatim}

We make this precise. Write ε\_n = (nπ/(2L))² · (sin(x\_n)/x\_n)² where x\_n =
nπ/(2L). Then

\begin{verbatim}
ε_n^{-s} = (2L/(nπ))^{2s} · (x_n/sin(x_n))^{2s}
\end{verbatim}

and

\begin{verbatim}
ζ_F(s)/(2L/π)^{2s} = Σ_{n=1}^{L-1} n^{-2s} · (x_n/sin(x_n))^{2s}
\end{verbatim}

\textbf{Error analysis.} The error has two sources:

\emph{(i) Tail truncation.} The missing terms Σ\_\{n≥L\} n\^{}\{-2s\} contribute
O(L\^{}\{1−2s\}) for Re(s) \textgreater{} 1/2.

\emph{(ii) Sine correction.} For each mode, (x/sin(x))\^{}\{2s\} = 1 + sx²/3 +
O(x⁴). Summing the leading correction:

\begin{verbatim}
Σ_{n=1}^{L-1} n^{-2s} · s·(nπ/(2L))²/3 = (sπ²/(12L²)) · Σ_{n=1}^{L-1} n^{2−2s}
\end{verbatim}

For s \textgreater{} 3/2 this sum converges and the correction is O(1/L²). For
1/2 \textless{} s ≤ 3/2, the sum grows as L\^{}\{3−2s\} and the correction is
O(L\^{}\{1−2s\}) = O(1/L) when s = 1.

The combined rate is \(O(L^{1-2\sigma})\) for \(1/2 < \sigma \leq 3/2\),
improving to \(O(1/L^2)\) for \(\sigma > 3/2\). ∎

\textbf{Remark.} This is a manifestation of Weyl's law {[}14{]}: for any
self-adjoint operator on a one-dimensional domain of length L with Dirichlet
boundary conditions, the spectral zeta function reproduces ζ(2s) through the
asymptotic eigenvalue density at the spectral edge. The carry bridge operator is
a concrete arithmetic realization of this universal phenomenon.

\subsubsection{7.3 Numerical verification}\label{numerical-verification}

Verified numerically for s = 1, 2, 3, 4 (even zeta values) and s = 3/2, 5/2 (odd
zeta values ζ(3), ζ(5)):

\begin{longtable}[]{@{}lllll@{}}
\toprule\noalign{}
s & ζ(2s) & Ratio at L=100 & Ratio at L=500 & Ratio at L=1000 \\
\midrule\noalign{}
\endhead
\bottomrule\noalign{}
\endlastfoot
1 & π²/6 = 1.6449 & 1.6776 & 1.6515 & 1.6482 \\
2 & π⁴/90 = 1.0823 & 1.1266 & 1.0910 & 1.0867 \\
3/2 & ζ(3) = 1.2021 & 1.2291 & 1.2075 & 1.2048 \\
5/2 & ζ(5) = 1.0369 & 1.0542 & 1.0404 & 1.0387 \\
\end{longtable}

All ratios converge to the exact ζ(2s) values at rates consistent with Theorem 3
(\(O(L^{1-2\sigma})\) for \(\sigma \leq 3/2\), \(O(1/L^2)\) for
\(\sigma > 3/2\)).

\subsubsection{7.4 Universality and
non-universality}\label{universality-and-non-universality}

The ζ(2s) convergence is \emph{universal}: it depends only on the Dirichlet
boundary conditions and the length L, not on the specific transition
probabilities. Any operator with the same boundary structure produces the same
result.

By contrast, the conjectured limit R → −π is \emph{not} universal --- it depends
on the binary carry structure, the Fejér kernel correlations, and the D-odd
constraint (Section 10).

\begin{center}\rule{0.5\linewidth}{0.5pt}\end{center}

\newpage

\subsection{8. The Bernoulli Hierarchy and the Gaussian
Integers}\label{the-bernoulli-hierarchy-and-the-gaussian-integers}

\subsubsection{8.1 The Bernoulli hierarchy}\label{the-bernoulli-hierarchy}

The carry chain in binary multiplication encodes Bernoulli numbers at successive
levels of its spectral decomposition:

\begin{longtable}[]{@{}llll@{}}
\toprule\noalign{}
Level & Carry quantity & Value & Bernoulli \\
\midrule\noalign{}
\endhead
\bottomrule\noalign{}
\endlastfoot
0 & Stationary distribution & 1 & B₀ = 1 \\
1 & Spectral gap λ₂ & 1/2 & \textbar B₁\textbar{} = 1/2 \\
2 & Sub-sector ratio α & 1/6 & B₂ = 1/6 \\
3 & Value cancellation σ₀₀\^{}\{c₂=1,c₃=1\} & 0 & B₃ = 0 \\
\end{longtable}

\textbf{Level 0} is trivial: the carry chain has a unique stationary
distribution. \textbf{Level 1} is the Diaconis--Fulman theorem {[}1{]}: the
spectral gap of the carry transfer operator in base b is 1/b, giving 1/2 =
\textbar B₁\textbar{} for binary. \textbf{Level 2} is new. The sub-sector ratio
α measures the carry population at position D−2, which is directly linked to the
boundary layer decomposition of {[}E, §8.5{]}: the c\_\{D−2\} = 1 sub-population
is precisely the ``top'' cascade contribution that is blocked in sector (1,0).

\textbf{Definition 5 (Sub-sector ratio).} Define α(K) =
σ₀₀\textsuperscript{\{c=1\}(K)/σ₀₀}\{c=0\}(K), where σ₀₀\^{}\{c=0\} and
σ₀₀\^{}\{c=1\} are the sector-(0,0) partition function restricted to pairs with
c\_\{D−2\} = 0 and c\_\{D−2\} = 1, respectively.

\textbf{Conjecture 2 (Bernoulli sub-sector ratio).} α(K) → 1/6 = B₂ as K → ∞.

\emph{The evidence is preliminary: the overshoot at K ≥ 20 means the conjecture
rests on Richardson extrapolation (3.3 digits) rather than direct convergence.
Additional data at higher K would strengthen the case.}

\emph{Numerical evidence.} Exact enumeration gives:

\begin{longtable}[]{@{}llll@{}}
\toprule\noalign{}
K & α(K) & \(\lvert \alpha - 1/6 \rvert\) & Sig. digits \\
\midrule\noalign{}
\endhead
\bottomrule\noalign{}
\endlastfoot
7 & 0.15177 & 0.0149 & 1.2 \\
10 & 0.15958 & 0.0071 & 1.5 \\
13 & 0.16378 & 0.0029 & 1.9 \\
16 & 0.16575 & 0.0009 & 2.4 \\
18 & 0.16652 & 0.0001 & 2.7 \\
20 & 0.16880 & 0.0021 & 2.0 \\
21 & 0.16922 & 0.0026 & 2.6 \\
\end{longtable}

The convergence is non-monotonic, with a two-phase structure: approach from
below through K = 18, then overshoot at K ≥ 20. The overshoot is present at K =
21 (α = 0.1692, gap = 0.0026). Richardson extrapolation with the ansatz α(K) =
α(∞) + a₁K²/2\^{}K + \ldots{} gives α(∞) ≈ 0.1667 (3.3 digits), consistent with
1/6. The correction structure is K²/2\^{}K rather than pure 1/2\^{}K. The
non-monotonic convergence is analogous to the oscillatory behaviour of R\_recon
in Paper E (§8.3), likely reflecting a Gibbs-type phenomenon in the spectral
truncation.

\subsubsection{8.2 The Euler--Maclaurin connection
(speculative)}\label{the-eulermaclaurin-connection-speculative}

\emph{This connection is a structural analogy rather than a formal derivation;
the numerical coincidence is striking but unproved.}

The carry chain computes ⌊(conv + carry)/2⌋ at each position --- a discrete
summation with rounding. The Euler--Maclaurin formula for the discrepancy
between a discrete sum and its integral is

\begin{verbatim}
Σf − ∫f = B₁·Δf + B₂/2!·Δf' + B₄/4!·Δf''' + ⋯
\end{verbatim}

where B\_n are Bernoulli numbers. The carry chain's spectral decomposition
mirrors this hierarchy: - B₁ = −1/2 controls the first-order correction
(exponential convergence rate = spectral gap); - B₂ = 1/6 controls the
second-order correction (the sub-sector ratio α, which measures how much of the
carry at D − 2 is ``rescued'' from the c = 1 sub-population).

This connection is structural rather than a formal derivation: the carry chain
is a nonlinear discrete dynamical system, not a trapezoidal sum. Nonetheless,
the numerical evidence strongly supports the identification.

\textbf{Theorem 6 (Cascade Rigidity).} For D-odd sector-(0,0) pairs with
\(c_2 \in \lbrace{}0,1\rbrace{}\):

\begin{enumerate}
\def\labelenumi{(\roman{enumi})}
\item
  \(\text{carries}[D-1] = \text{carries}[D] = 0\).
\item
  If \(c_2 = \text{carries}[D-2] \geq 1\): \(M = D-2\) and
  \(\text{val} = c_3 - 1\).
\item
  If \(c_2 = 0\) and \(c_3 = \text{carries}[D-3] \geq 1\): \(M = D-3\) and
  \(\text{val} = c_4 - 1\).
\item
  If \(c_2 = c_3 = 0\): \(M\) depends on the full carry chain.
\end{enumerate}

\emph{Proof.} In sector (0,0), \(x_{K-2} = y_{K-2} = 0\) and
\(x_{K-1} = y_{K-1} = 1\). The convolution at position \(D-2\) is
\(\text{conv}_{D-2} = x_{K-2} y_{K-1} + x_{K-1} y_{K-2} = 0\). Then
\(\text{carries}[D-1] = \lfloor(0 + c_2)/2\rfloor = 0\) for
\(c_2 \in \lbrace{}0,1\rbrace{}\), and
\(\text{carries}[D] = \lfloor(1 + 0)/2\rfloor = 0\). Scanning from \(D\)
downward, the first nonzero carry is at \(D-2\) when \(c_2 \geq 1\), giving
\(M = D-2\) and \(\text{val} = \text{carries}[D-3] - 1 = c_3 - 1\). When
\(c_2 = 0\), the scan continues to \(D-3\), etc. □

The theorem extends recursively: at each level \(n \geq 2\), if
\(c_n = \text{carries}[D-n] \geq 1\) while \(c_2 = \cdots = c_{n-1} = 0\), then
\(M = D-n\) and \(\text{val} = c_{n+1} - 1\). This recursive extension is
verified numerically through \(n = 6\) at \(K \leq 13\); a general inductive
proof for all \(n\) is not given.

\textbf{Corollary (Level 3: B₃ = 0).} For \(c_2 \geq 1\), Theorem 6(ii) gives
\(\text{val} = c_3 - 1\). When \(c_3 = 1\), every pair contributes
\(\text{val} = 0\), so \(\sigma_{00}^{c_2 \geq 1,\, c_3 = 1} = 0\) exactly. This
is the carry-chain incarnation of \(B_3 = 0\): the corresponding term in the
Euler--Maclaurin formula vanishes. Verified for all \(K \leq 13\)
(\(n = 1{,}789{,}942\) pairs at \(K = 13\)).

The recursive extension produces the same vanishing at every depth:
\(\sigma_{00}^{c_2 = 0,\, c_3 \geq 1,\, c_4 = 1} = 0\), and so on. At each
level, only the ``all-zeros'' sub-population \((c_2 = c_3 = \cdots = c_n = 0)\)
carries non-trivial cascade dynamics.

\textbf{Level 4: open.} The non-trivial spectral content telescopes into the
\((c_2, c_3, c_4) = (0,0,0)\) and \((0,0,1)\) cells. Seven candidate
partition-function ratios, the asymptotic correction structure of \(\alpha(K)\),
and the spectral gap of the near-boundary carry transfer matrix were tested
through \(K = 13\); none produces \(|B_4| = 1/30\). Identifying the correct
observable remains open.

\subsubsection{8.2a Stopping-time
decomposition}\label{a-stopping-time-decomposition}

Theorem 6, combined with the unit leading carry (Theorem 7, §10.5), yields an
exact stopping-time decomposition of the sector ratio:

\[\sigma_{ab} = \sum_{\tau=2}^{D-2} n_{ab}(\tau) \cdot E[\text{val} \mid \tau,\, ab]\]

where \(\tau = D - M\) is the depth of the first nonzero carry from the MSB. The
Cascade Rigidity Theorem makes this decomposition rigorous: at each \(\tau\),
val \(= c_{\tau+1} - 1\) where \(c_{\tau+1} = \text{carries}[D-\tau-1]\).

\textbf{Structural exclusion.} In sector \((1,0)\), \(\text{conv}_{D-2} = 1\)
(Proposition 4), so \(\text{carries}[D-2] = 0\) and
\(P(\tau = 2 \mid \text{sector } 10) = 0\) for all \(K\). Sector \((0,0)\) has
\(\text{conv}_{D-2} = 0\), allowing \(\tau = 2\).

\textbf{Universality.} Exact enumeration confirms that both \(P(\tau \mid ab)\)
and \(E[\text{val} \mid \tau, ab]\) converge to \(K\)-independent constants.

\textbf{Table 1.} Stopping-time decomposition constants at \(K = 21\) (from E45
exact enumeration, \(3.4 \times 10^{11}\) D-odd pairs in the \(X \leq Y\)
half-plane); cross-K stability from \(K = 19\) to \(K = 21\) confirms
convergence to 5--6 significant digits.

\begin{longtable}[]{@{}
  >{\raggedright\arraybackslash}p{(\columnwidth - 10\tabcolsep) * \real{0.0870}}
  >{\raggedleft\arraybackslash}p{(\columnwidth - 10\tabcolsep) * \real{0.1848}}
  >{\raggedleft\arraybackslash}p{(\columnwidth - 10\tabcolsep) * \real{0.1848}}
  >{\raggedleft\arraybackslash}p{(\columnwidth - 10\tabcolsep) * \real{0.1957}}
  >{\raggedleft\arraybackslash}p{(\columnwidth - 10\tabcolsep) * \real{0.1957}}
  >{\raggedleft\arraybackslash}p{(\columnwidth - 10\tabcolsep) * \real{0.1522}}@{}}
\toprule\noalign{}
\begin{minipage}[b]{\linewidth}\raggedright
\(\tau\)
\end{minipage} & \begin{minipage}[b]{\linewidth}\raggedleft
\(P(\tau\mid 00)\)
\end{minipage} & \begin{minipage}[b]{\linewidth}\raggedleft
\(P(\tau\mid 10)\)
\end{minipage} & \begin{minipage}[b]{\linewidth}\raggedleft
\(E[v\mid\tau,00]\)
\end{minipage} & \begin{minipage}[b]{\linewidth}\raggedleft
\(E[v\mid\tau,10]\)
\end{minipage} & \begin{minipage}[b]{\linewidth}\raggedleft
\(\Delta(\tau)\)
\end{minipage} \\
\midrule\noalign{}
\endhead
\bottomrule\noalign{}
\endlastfoot
2 & 0.54070 & 0.00000 & −0.00812 & --- & --- \\
3 & 0.13713 & 0.45726 & −0.11761 & +0.26971 & +0.38732 \\
4 & 0.10924 & 0.22875 & −0.04089 & +0.31022 & +0.35111 \\
5 & 0.08557 & 0.14398 & −0.07338 & +0.33738 & +0.41076 \\
6 & 0.05007 & 0.07849 & −0.00867 & +0.30657 & +0.31525 \\
7 & 0.03329 & 0.04308 & −0.01270 & +0.32311 & +0.33580 \\
8 & 0.01862 & 0.02313 & +0.04186 & +0.33033 & +0.28848 \\
9 & 0.01169 & 0.01231 & +0.01742 & +0.33307 & +0.31565 \\
\end{longtable}

The sector asymmetry \(\Delta(\tau) = E[v\mid\tau,10] - E[v\mid\tau,00]\) is
stable to \(< 0.00001\) across \(K = 19\)--\(21\) for \(\tau \leq 5\),
confirming genuine universality. Richardson extrapolation (assuming convergence
rate \(\rho = 1/2\)) yields 7-digit estimates of the limiting constants.

\textbf{Identification status.} PSLQ searches against
\(\lbrace{}1, \pi, \ln 2, \ln 3\rbrace{}\) on the Richardson-extrapolated limits
produce candidate relations with residuals in the range \(10^{-6}\)--\(10^{-7}\)
but no confident closed forms. The precision is sufficient to confirm
universality but not to prove algebraic identities. Extending the stopping-time
analysis to \(K \geq 24\) (computationally feasible with the E45 C
implementation) would provide the 10+ digits needed for reliable PSLQ
identification.

\subsubsection{8.2b Failure of the per-position spectral
limit}\label{b-failure-of-the-per-position-spectral-limit}

A natural approach to the sector ratio is to analyze the per-position carry
transition matrices \(T(d)\) directly. Under D-odd conditioning, the carry chain
at each position \(d\) induces a stochastic matrix on carry states
\(\lbrace{}0, 1, 2, \ldots\rbrace{}\). In the Markov (independent-convolution)
model, the bulk eigenvalue is \(\lambda_2 = 1/2\) at every position; the Linear
Mix Hypothesis {[}E{]} would predict
\(\lambda_2 = (1-A^{\ast})/2 + A^{\ast}/2 \approx 0.69\).

Empirical computation of the bulk-averaged transition matrix for
\(K = 5\)--\(12\) (experiment P1\_05) reveals a qualitatively different picture:

\begin{longtable}[]{@{}lcc@{}}
\toprule\noalign{}
K & \(\lambda_2(\text{bulk})\) & \(A_{\text{eff}}\) \\
\midrule\noalign{}
\endhead
\bottomrule\noalign{}
\endlastfoot
6 & 0.690 & −3.25 \\
8 & 0.756 & −0.88 \\
10 & 0.810 & −0.40 \\
12 & 0.858 & −0.19 \\
\end{longtable}

The second eigenvalue \(\lambda_2\) grows monotonically toward 1, and the
effective mixing parameter \(A_{\text{eff}}\) extracted from the LMH formula
diverges to \(-\infty\) instead of converging to \(A^{\ast} \approx 1.38\). The
per-position transition matrices become \emph{more rigid} (closer to identity)
as \(K\) increases --- the D-odd conditioning progressively constrains the local
transitions, which is the dynamical mechanism behind Theorem 6 (Cascade
Rigidity).

\textbf{Conclusion.} The Linear Mix Hypothesis cannot hold at the per-position
level. The LMH must describe a \emph{macroscopic} property of the full resolvent

\[(I - \mathcal{K}_{\text{eff}})^{-1}\]

, not a uniform modification of each per-position factor. This is consistent
with the resolvent universality mechanism identified in {[}E, §8.3{]}: the
spectral resolvent achieves \(-\pi\) through collective mode weighting, not
through any single position's transition structure. The stopping-time
decomposition (§8.2a) provides the correct analytical framework precisely
because it bypasses the per-position transition matrices entirely.

\textbf{Boundary-layer concentration.} Approximately 95\% of \(R(K)\) comes from
\(\tau \leq 7\). The cumulative ratio through \(\tau = 7\) at \(K = 14\) is
\(-2.66\), close to the full \(R(14) = -2.62\). The cascade dynamics is
overwhelmingly a boundary-layer phenomenon.

\textbf{Even-odd alternation.} The successive ratios
\(P(\tau+1 \mid 00)/P(\tau \mid 00)\) exhibit a clear even-odd oscillation
converging toward \(1/2\):
\(\lbrace{}0.254, 0.796, 0.783, 0.584, 0.664, 0.557, 0.626, \ldots\rbrace{}\).
This is consistent with the Diaconis--Fulman spectral gap \(\rho = 1/2\) as the
asymptotic rate, modulated by the carry chain's non-stationary structure near
the boundary.

\subsubsection{8.3 The ζ(2) product
(preliminary)}\label{the-ux3b62-product-preliminary}

\emph{This claim is conditional on two unproved conjectures (α → 1/6 and R →
−π); the ``new proof of the Basel problem'' is speculative until both are
established.}

The combination of levels 1 and 2 with the conjectural R = −π gives

\begin{verbatim}
α · R² = (1/6) · π² = ζ(2)
\end{verbatim}

If both α = 1/6 (Conjecture 2) and R = −π (Conjecture 1) can be established
analytically from carry arithmetic, this would yield a new proof of the Basel
problem ζ(2) = π²/6 --- one that begins with integer multiplication rather than
Fourier analysis or contour integration.

\emph{Numerical verification.} The product α(K) · R(K)² converges:

\begin{longtable}[]{@{}lll@{}}
\toprule\noalign{}
K & α · R² & \(\lvert \alpha \cdot R^2 - \zeta(2) \rvert\) \\
\midrule\noalign{}
\endhead
\bottomrule\noalign{}
\endlastfoot
11 & 0.405 & 1.24 \\
15 & 1.218 & 0.43 \\
18 & 1.583 & 0.062 \\
20 & 1.648 & 0.003 \\
21 & 1.662 & 0.017 \\
\end{longtable}

The product α · R² overshoots ζ(2) at K = 21 because α overshoots 1/6 while R
has not yet reached −π. The near-exact agreement at K = 20 (gap = 0.003) is
partially accidental: it reflects a cancellation between the α overshoot and the
R undershoot at that particular K.

\subsubsection{8.4 Base specificity}\label{base-specificity}

The sub-sector ratio α = 1/6 is \textbf{not universal} --- it is specific to
base 2. Monte Carlo experiments in bases 3, 5, and 7 show no convergence to 1/6
or to any other Bernoulli number. In base 3, α is wildly unstable across K
values; in base 5, the sector ratio itself appears rational (R₅ → 5/4) with α₅
far from 1/6. This base specificity makes the binary Bernoulli hierarchy more,
not less, remarkable: it reflects a deep structural property of base-2
arithmetic, potentially connected to the ramification of 2 in ℤ{[}i{]} (Section
8.6). See {[}G{]} for the proof that base 2 is the unique base where the
D-parity boundary is a straight line in angular coordinates.

\subsubsection{8.5 The Euler product
perspective}\label{the-euler-product-perspective}

Fulman's multiplicativity property K\_a · K\_b = K\_\{ab\} for carry transfer
operators {[}3{]} implies that the Markov spectral sum in base b factors as

\begin{verbatim}
S_b(0) = b²/(b² − 1)
\end{verbatim}

(verified numerically for b = 2, 3, 5, 7, 10). For b = 2 this gives S₂(0) = 4/3.
(Note: this spectral zeta sum S₂(0) is a different quantity from the Markov
sector ratio R\_Markov = +2/3 of Theorem 2; both derive from
independent-convolution statistics but measure different observables.) Taking
the product over all primes:

\begin{verbatim}
∏_{p prime} S_p(0) = ∏_p p²/(p² − 1) = ζ(2)
\end{verbatim}

This is a restatement of the classical Euler product, but viewed through carry
arithmetic: each prime p contributes its Markov spectral sum factor to ζ(2). The
non-Markov (exact) sector ratio introduces L(1, χ₄) = π/4 as the additional
factor, connecting to the Dedekind zeta function (Section 8.6).

\subsubsection{8.6 The Dedekind zeta function of ℚ(i)
(speculative)}\label{the-dedekind-zeta-function-of-ux211ai-speculative}

\emph{This is a heuristic connection: the factorization is classical;
attributing both factors to carry arithmetic is conjectural and depends on
Conjecture 1.}

The Gaussian integers ℤ{[}i{]} have Dedekind zeta function

\begin{verbatim}
ζ_{ℚ(i)}(s) = ζ(s) · L(s, χ₄)
\end{verbatim}

The carry bridge produces: - ζ(2s) from the spectral zeta function (Theorem 3),
via the Markov eigenvalue density - L(1, χ₄) = π/4 from the sector ratio
(Conjecture 1), via the effective χ₄ model (§6.2) and resolvent universality
{[}E{]}

Both factors of ζ\_\{ℚ(i)\} thus emerge from the same arithmetic object --- the
carry bridge of binary multiplication. The prime 2 is the unique rational prime
that \emph{ramifies} in ℤ{[}i{]}, satisfying 2 = −i(1+i)². This ramification
forces the character χ₄ into the spectral decomposition, explaining at the
algebraic level why base 2 is unique {[}G{]}: for b ≥ 3, the sector ratio
appears to be rational (e.g., R = 5/4 for base 5), and the D-parity boundary is
curved in angular coordinates, removing the geometric mechanism that produces
L(1, χ₄).

\begin{center}\rule{0.5\linewidth}{0.5pt}\end{center}

\newpage

\subsection{9. Numerical Evidence}\label{numerical-evidence}

\subsubsection{9.1 Exact enumeration}\label{exact-enumeration}

The sector ratio R(K) is computed by exhaustive enumeration of all 4\^{}\{K−1\}
digit pairs. Enumeration through K = 21 covers 5.5 × 10¹¹ pairs for K = 21 (C
implementation with 10-thread parallelism, \textasciitilde2.7 hours) {[}5{]}.

\begin{longtable}[]{@{}llll@{}}
\toprule\noalign{}
K & R(K) & \(\lvert R(K) + \pi \rvert\) & Sig. digits \\
\midrule\noalign{}
\endhead
\bottomrule\noalign{}
\endlastfoot
7 & −0.092 & 3.050 & --- \\
9 & −0.726 & 2.416 & --- \\
11 & −1.554 & 1.588 & --- \\
13 & −2.339 & 0.803 & 0.1 \\
15 & −2.824 & 0.317 & 0.5 \\
17 & −3.035 & 0.107 & 1.0 \\
19 & −3.110 & 0.032 & 1.5 \\
20 & −3.125 & 0.017 & 1.8 \\
21 & −3.133 & 0.008 & 2.1 \\
\end{longtable}

All values from the cascade valuation (Definition 1), cross-validated against
{[}E, Table 1{]}.

\subsubsection{9.2 Richardson extrapolation}\label{richardson-extrapolation}

The raw convergence R(K) → −π is slow: at K = 21 we have 2.1 significant digits.
Richardson extrapolation dramatically improves precision.

\textbf{Ansatz 1 (polynomial).} A natural first guess is R(K) = R(∞) + a₁/K² +
a₂/K⁴ + ⋯, motivated by the O(1/L²) finite-size correction in the Möbius sum
S(A, L) (Proposition 2). However, the true convergence rate is exponential
(§9.3), and the polynomial Neville table diverges in practice.

\textbf{Ansatz 2 (exponential).} The empirical gap ratios g(K)/g(K−1) converge
monotonically to 1/2 (§9.3, Table 1 of {[}E{]}), consistent with

\begin{verbatim}
R(K) = R(∞) + c₁ ρ^K + c₂ ρ^{2K} + c₃ ρ^{3K} + ⋯,    ρ = 1/2.
\end{verbatim}

The Neville--Aitken Richardson table with step-size variable h = ρ\^{}K
eliminates successive error terms.

\textbf{Proposition 3 (Richardson extrapolation).} Applying the Neville--Aitken
deferred approach to the limit {[}9{]} to the sequence R(K) for K = 7, 9, 11,
\ldots, 21 with Ansatz 2 (h = (1/2)\^{}K) yields

\begin{verbatim}
R_extrap = −3.1419… (4.0 significant digits)
\end{verbatim}

using exact R(K) for K ≤ 13 and K ≥ 19, with E44 approximations for K = 14--18.
Using only exact values (K ≤ 13, K ≥ 19) the estimate improves to 4.4
significant digits. Leave-one-out cross-validation confirms stability: removing
any single K ∈ \{7, \ldots, 17\} changes the estimate by less than 10⁻⁴.

\subsubsection{9.3 Convergence rate}\label{convergence-rate}

The error ε(K) = \textbar R(K) + π\textbar{} satisfies

\begin{verbatim}
ε(K) = O(K² · 2^{−K})
\end{verbatim}

for K ≥ 11. The consecutive gap ratio g(K)/g(K−1) forms a monotone sequence
converging to 1/2 from above: 1.000, 0.954, \ldots, 0.522, 0.491 at K = 21 (see
{[}E, Table 1{]}). The gap ratio crossing 2 at K = 21 confirms the asymptotic
regime. This exponential rate is consistent with the Diaconis--Fulman spectral
radius ρ = 1/2 governing the carry chain {[}4{]}.

\emph{Remark.} A polynomial fit ε ∼ C/K² provides a reasonable approximation for
moderate K but misses the asymptotic structure. The exponential characterization
ε ∼ C · (1/2)\^{}K captures the true rate. The Richardson extrapolation of §9.2
remains effective because at moderate K the polynomial and exponential forms are
numerically similar.

\subsubsection{9.4 Verification of the closed
form}\label{verification-of-the-closed-form}

Theorem 1 is verified numerically at 50-digit precision using
arbitrary-precision arithmetic. The algebraic identity 4r/(γ(1−r²)) = 2β/(1+2β)
holds to full precision for all tested values of β. The formula S(β*, L) → −π is
confirmed at L = 1000 with \textbar S + π\textbar{} \textless{} 4 × 10⁻⁵,
consistent with the O(1/L²) rate.

\begin{center}\rule{0.5\linewidth}{0.5pt}\end{center}

\newpage

\subsection{10. The Gap: From S(A) to R(K)}\label{the-gap-from-sa-to-rk}

\subsubsection{10.1 What is proved}\label{what-is-proved}

Theorem 1 establishes that the abstract alternating sum S(A) = 2(1−A)/(3−2A) is
a closed-form Möbius transformation. Setting S = −π yields A* = (3π+2)/(2(π+1)).
The spectral zeta function ζ\_F(s) → ζ(2s) is proved (Theorem 3). The Dirichlet
character χ₄(n) = sin(nπ/2) emerges as the effective spectral weight (§6.2).

\subsubsection{10.2 The gap and the universality
phenomenon}\label{the-gap-and-the-universality-phenomenon}

The gap between the abstract sum S(A) and the physical sector ratio R(K) is the
central open problem. Three observations constrain its nature:

\textbf{Observation 1 (Distinct convergence rates).} The abstract sum converges
polynomially: S(A, L) − S(A, ∞) = O(1/L²). The physical sector ratio converges
exponentially: R(K) − R(∞) = O(K²/2\^{}K).

\textbf{Observation 2 (A\_fit does not converge).} Defining A\_fit(K) by the
equation S(A\_fit, 2K−1) = R(K), the sequence A\_fit(K) does not converge to A*.

\textbf{Observation 3 (Carry projection is exact Markov).} Projecting the exact
(non-Markov) transfer operator onto carry degrees of freedom recovers the Markov
transfer matrix exactly. The parameter A is therefore \emph{not} a perturbative
eigenvalue shift.

These observations rule out deriving A* from the Fejér kernel correlation
structure. Nonetheless, both S(A, L) and R(K) converge numerically to the same
value −π --- a \emph{universality} phenomenon {[}P2{]} that we state as a
conjecture:

\textbf{Conjecture 3 (Universality).} Let
\(S(A^{\ast}, L) = 2(1-A^{\ast})/(3-2A^{\ast})\) be the abstract Möbius sum
(Theorem 1), and let \(R(K) = \sigma_{10}(K)/\sigma_{00}(K)\) be the cascade
sector ratio (Definition 2). Then

\[\lim_{L \to \infty} S(A^{\ast}, L) \;=\; \lim_{K \to \infty} R(K) \;=\; -\pi,\]

\emph{where the left equality is proved (Theorem 1) and the right is Conjecture
1. The two systems share different microscopic dynamics (polynomial vs
exponential finite-size corrections) but the same macroscopic limit.}

The common mechanism is the Dirichlet character χ₄ and the L-function L(1, χ₄) =
π/4. {[}E{]} identifies \textbf{resolvent universality} as the specific
mechanism: the spectral resolvent

\[(I - \mathcal{K}_{\mathrm{eff}})^{-1}\]

, applied to the true per-sector cascade profiles, produces a macroscopic sum
converging to \(-\pi\) through collective mode weighting. The left limit is an
exact algebraic evaluation; the right is supported by 4.0-digit exponential
Richardson extrapolation through K = 21 (§9).

\subsubsection{10.3 Structural decomposition}\label{structural-decomposition}

\textbf{Proposition 4 (Top-carry constraint).} For a D-odd product n = p·q with
D = 2K−1, the convolution at position D−2 equals exactly conv\_\{D−2\} = a + c,
where a = p\_\{K−2\} and c = q\_\{K−2\} are the sector bits.

\emph{Proof.} conv\_\{D−2\} = p\_\{K−2\} · q\_\{K−1\} + p\_\{K−1\} · q\_\{K−2\}
= a · 1 + 1 · c = a + c.~□

\textbf{Corollary.} The D-odd condition requires carries{[}D−1{]} = 0, hence (a
+ c + carries{[}D−2{]}) \textless{} 2. For sector (1,0): a + c = 1, so
carries{[}D−2{]} = 0, forcing M ≤ D−3. For sector (0,0): carries{[}D−2{]} ∈
\{0,1\}, yielding two sub-populations. This gives σ\_\{00\} =
σ\_\{00\}\^{}\{top\} + σ\_\{00\}\^{}\{low\}, while σ\_\{10\} =
σ\_\{10\}\^{}\{low\} has no top contribution. {[}E, §8.5{]} independently
establishes this structural exclusion: P(j=2\textbar sector 10) = 0, confirmed
by backward dynamic programming ({[}E{]}).

\textbf{Proposition 5 (Exact symmetry).} σ\_\{10\}(K) = σ\_\{01\}(K) for all K.

\emph{Proof.} The substitution (p,q) ↦ (q,p) maps sector (1,0) to (0,1),
preserving all carries. □

\textbf{Theorem 4 (N₁₁ = 0).} For all K ≥ 2, no D-odd pair has sector (1,1).

\emph{Proof.} In sector (1,1), conv\_\{D−2\} = 2. The D-odd condition requires
carries{[}D−1{]} = 0, but (2 + carries{[}D−2{]})/2 ≥ 1, contradiction. □

\textbf{Theorem 5 (Closed-form count ratio).} In the continuous limit K → ∞,

\[N_{10}/N_{00} \to \frac{2\ln(4/3) - 1/2}{2\ln(9/8)} = 0.31992784225\ldots\]

with correction O(1/2\^{}K).

\emph{Proof.} Write X = 2\^{}\{K−1\}(1 + u), Y = 2\^{}\{K−1\}(1 + v) with u, v ∈
{[}0, 1). The D-odd condition (1 + u)(1 + v) \textless{} 2 defines a region in
(u,v). Sectors are determined by ⌊2u⌋ and ⌊2v⌋: area f₀₀ = 2 ln(9/8), f₁₀ = 2
ln(4/3) − 1/2, f₁₁ = 0. Verified to 8 significant digits against exact
enumeration (K = 3--21). □

\textbf{Corollary (π isolation).} Since R = (⟨val⟩₁₀/⟨val⟩₀₀) × (N₁₀/N₀₀) and
N₁₀/N₀₀ involves only ln 2 and ln 3, the entire factor π enters through the
carry weight ratio:

\[\langle\text{val}\rangle_{10}/\langle\text{val}\rangle_{00} \to -\pi \cdot \frac{2\ln(9/8)}{2\ln(4/3) - 1/2} = -9.8197\ldots\]

Any proof must explain how the sector bit perturbation modifies the
\emph{distribution} of c\_\{M−1\} to produce this ratio. The Dirichlet series
approach (character decomposition by n mod 4) does not cleanly factor R, since
\textasciitilde75\% of D-odd products are even.

\emph{Remark.} {[}E, §9.2{]} falsifies the hypothesis that the \emph{cascade
count} ratio n₀₀/n₁₀ converges to π; the limit is algebraic (≈ 25/8). Combined
with Theorem 5, this confirms that π is entirely \emph{dynamical}: it resides in
the carry weight functional, not in any geometric count.

\subsubsection{10.4 The carry weight
mechanism}\label{the-carry-weight-mechanism}

The total variation distance TV(j) between the carry distributions in sectors
(1,0) and (0,0) reveals a boundary layer with two peaks: TV ≈ 0.16 at j = K−2
(sector bit perturbation) and TV ≈ 0.54 at j = D−2 (top structural constraint),
with a plateau TV ≈ 0.06 in between.

Four steps link the sector bit to the carry weight val = c\_\{M−1\} − 1:

\begin{enumerate}
\def\labelenumi{\arabic{enumi}.}
\item
  \textbf{Convolution perturbation (position K−2).} The sector bit a = 1 adds
  q\_\{j−K+2\} to conv\_j for j ∈ {[}K−2, 2K−3{]}, generating additional carries
  in the bulk.
\item
  \textbf{Carry propagation.} The extra carries propagate upward through
  \textasciitilde K positions; the mean carry at each position is
  \textasciitilde0.35 higher in sector (1,0) {[}F{]}.
\item
  \textbf{Top constraint (position D−2).} The D-odd condition forces
  carries{[}D−2{]} = 0 in sector (1,0), eliminating the ``top'' population and
  shifting M to D−3 or lower.
\item
  \textbf{Val distribution shift.} The higher carry levels in sector (1,0),
  combined with the constrained top, produce 30\% with c\_\{M−1\} = 2 (val = +1)
  vs 5.5\% in sector (0,0)\emph{\{low\}, driving σ}\{10\} \textgreater{} 0 while
  σ\_\{00\} \textless{} 0.
\end{enumerate}

\textbf{Spectral connection.} The sector bit perturbation is a ramp (constant
⟨Δconv\_j⟩ = 1/2) over half the chain. In the bridge eigenmode basis sin(nπj/L),
the ramp Fourier coefficients Δv\_n = {[}cos(nπ/2) − cos(nπ){]}/(nπ) combined
with the source-point factor sin(nπ/2) satisfy:

\[\sin(n\pi/2) \cdot [\cos(n\pi/2) - \cos(n\pi)]
= \sin(n\pi/2)\ \text{if } n \text{ is odd},\ \text{and } 0\ \text{if } n \text{ is even}.\]

The even-mode contributions cancel exactly; the odd-mode sum becomes Σ
sin(nπ/2)/n = π/4 = L(1, χ₄). This explains conceptually how the Leibniz series
enters the problem.

\textbf{Remaining normalization gap.} With the Diaconis--Fulman spectral gap λ =
1/2 taken as constant, the spectral sum evaluates to 1/2, not −π. The
discrepancy factor −2π must originate from: (i) the mode-dependent eigenvalue
spectrum of the non-Markovian carry chain; (ii) the non-local observation
functional val = c\_\{M−1\} − 1 evaluated at the random stopping position M−1;
and (iii) the redistribution of the M = D−2 population (\textasciitilde55\% of
sector (0,0)) into M ≤ D−3 in sector (1,0). All three concern carry weights,
consistent with the π-isolation corollary. The analytical resolution is
identified in {[}E, §9{]} as the \textbf{Holte bridge / Doob h-transform}
matching between the bulk carry distribution (conditioned on D-odd) and the
boundary layer exclusions.

\begin{center}\rule{0.5\linewidth}{0.5pt}\end{center}

\subsubsection{10.5 Analytical evidence: continuum limit and exact symbolic
forms}\label{analytical-evidence-continuum-limit-and-exact-symbolic-forms}

\textbf{Carry recursion from MSB.} Define carry\([d]\) = carry at position
\(D-1-d\) from below, with recursion
carry\([d] = 2 \cdot\)carry\([d-1] + b_d - C_d\) and carry\([0] = 0\).

\textbf{Theorem 7 (Unit leading carry).} Let \(d_0\) be the smallest
\(d \geq 1\) with carry\([d] \geq 1\). Then carry\([d_0] = 1\) exactly.

\emph{Proof.} Since carry\([d_0-1] = 0\), carry\([d_0] = b_{d_0} - C_{d_0}\).
For carry \(\geq 1\) with \(b_{d_0} \in \lbrace{}0,1\rbrace{}\) and
\(C_{d_0} \geq 0\): necessarily \(C_{d_0} = 0\) and \(b_{d_0} = 1\). □

\textbf{Corollary (Universal val formula).}
\(\text{val} = 1 + b_{d_0+1} - C_{d_0+1}\) for all depths \(d_0\) and both
sectors.

The continuum function
\(\text{val}_\infty(u,v) = \lim_{K \to \infty} \text{val}(u,v;K)\) is a
well-defined step function taking values in \(\lbrace{}-1,0,+1\rbrace{}\),
giving sector sums
\(\sigma_{ac} = \int\!\!\int_{\text{sector } ac,\, \text{D-odd}} \text{val}_\infty(u,v)\, du\, dv\).

\textbf{Exact symbolic forms.} Each \(\sigma(d_0)\) admits an exact closed form
as a rational linear combination of
\(\lbrace{}1, \ln p : p \text{ prime}, p \leq 2^{d_0+2}\rbrace{}\):

\[\sigma_{00}(1) = \tfrac{1}{2} + \tfrac{9}{2}\ln 2 - 2\ln 3 - 3\ln 5 + \tfrac{7}{4}\ln 7\]

The individual coefficients grow as \(\sim 4^{d_0}\) while
\(\sigma(d_0) = O(2^{-d_0})\). After 7 depth levels, the cumulative \(\ln 2\)
coefficient reaches \textasciitilde1523, while the total
\(\sigma_{00} \approx -0.007\) --- catastrophic cancellation (ratio \(> 10^5\))
essential to the emergence of \(\pi\). At each depth \(d_0\), new primes in
\([2^{d_0+1}, 2^{d_0+2}]\) enter, so the infinite series involves \textbf{all
primes}. By Baker's theorem, no finite combination of logarithms of algebraic
numbers equals \(\pi\); the infinite series circumvents this obstruction.

\begin{longtable}[]{@{}llll@{}}
\toprule\noalign{}
\(d_{\max}\) & \(\sigma_{10}/\sigma_{00}\) & \(\lvert R + \pi \rvert\) &
digits \\
\midrule\noalign{}
\endhead
\bottomrule\noalign{}
\endlastfoot
5 & \(-2.695\) & \(4.5 \times 10^{-1}\) & 0.85 \\
10 & \(-3.098\) & \(4.3 \times 10^{-2}\) & 1.86 \\
12 & \(-3.129\) & \(1.2 \times 10^{-2}\) & 2.41 \\
14 & \(-3.139\) & \(2.1 \times 10^{-3}\) & 3.18 \\
15 & \(-3.14146\) & \(1.4 \times 10^{-4}\) & \textbf{4.36} \\
\end{longtable}

Beyond \(d_0 = 15\), catastrophic cancellation exceeds float64 precision.
Extended precision (60 digits; {[}E{]}) reaches \(d_0 = 18\) but reveals a
structural limitation: the depth-first carry-free enumeration captures only
17--33\% of the D-odd configuration space, causing the ratio to worsen beyond
\(d_0 = 15\). This is a method limitation, not a failure of the \(-\pi\)
conjecture.

\newpage

\subsection{11. Open Problems}\label{open-problems}

\begin{enumerate}
\def\labelenumi{\arabic{enumi}.}
\item
  Prove \(\sigma_{10}/\sigma_{00} = -\pi\) analytically. The decomposition
  \(R = R_0 + \Delta R\) (where \(R_0 = -3.9312\ldots\) involves only \(\ln 2\)
  and \(\ln 3\), and \(\Delta R = +0.7896\ldots\) carries the entire
  transcendental content) reduces the problem to showing that the cascade
  correction \(\Delta R = |R_0| - \pi\). Six equivalent formulations are
  catalogued in §10.5; the depth series \(\sum_{d=1}^\infty \sigma(d)\)
  converges to 4.36 significant digits at \(d_{\max} = 15\). The most promising
  analytical path is the \textbf{stopping-time series} (§8.2a): derive
  \(P(\tau \mid ab)\) and \(E[\text{val} \mid \tau, ab]\) analytically from the
  carry chain structure, then sum
  \(R = \sum_\tau [n_{10}(\tau) \cdot E[\text{val} \mid \tau, 10] + \cdots]\).
  The sector asymmetry
  \(\Delta(\tau) = E[v \mid \tau, 10] - E[v \mid \tau, 00]\), which converges to
  \(K\)-independent constants (Table 1, §8.2a), is the core quantity: it encodes
  the entire transcendental content of R. The per-position spectral approach
  (Doob h-transform of local transition matrices) is ruled out: §8.2b shows that
  the per-position eigenvalues grow with \(K\) instead of stabilizing, so the
  LMH must operate at the global resolvent level, not locally. Current status
  from the Step3 bridge: the \texttt{χ₄\ -\textgreater{}\ L(1,χ₄)} part is
  reduced to a scalar closure \texttt{R(∞)=C·L(1,χ₄)} with data-constrained
  \texttt{C∈{[}-4.0032,-3.9896{]}}, and the dominant decay rate is reduced to
  proving the analytic \texttt{1/2} bound for the conditioned chain. See {[}E,
  §9{]} for the updated closure map.
\item
  \textbf{Base universality.} For base b ≥ 3, the Markov sum S\_b(0) =
  b²/(b²−1). Is the exact sector ratio R\_b rational for all b ≥ 3? Exact
  enumeration for base 3 and base 5 would test this. The conjecture R₅ = 5/4 has
  preliminary numerical support.
\item
  \textbf{Higher Bernoulli levels.} The hierarchy B₀, \textbar B₁\textbar, B₂ is
  confirmed through Level 2. Level 3 (B₃ = 0) is proved by Theorem 6: the
  cascade rigidity forces \(\sigma_{00}^{c_2 \geq 1,\, c_3=1} = 0\) exactly, and
  the recursive extension produces the same vanishing at every depth. The Level
  4 signal \textbar B₄\textbar{} = 1/30 does not appear as a simple
  partition-function ratio, asymptotic correction coefficient, or spectral gap
  (all tested through K = 13). The non-trivial dynamics telescopes into the
  \((c_2, c_3, \ldots) = (0, 0, \ldots)\) sub-population at each level.
  Identifying the correct Level 4 observable remains open.
\item
  \textbf{Catalan's constant from carries.} The value L(2, χ₄) = G =
  0.9159\ldots{} (Catalan's constant) is the next L-function value after L(1,
  χ₄) = π/4. Can it be extracted from carry chain data?
\item
  \textbf{ζ(s) at odd arguments} (speculative). The Dedekind factorization ζ(s)
  = ζ\_\{ℚ(i)\}(s)/L(s, χ₄) gives ζ(s) at all arguments if L(s, χ₄) is
  accessible from carries for general s. The case ζ(3) = 1.20206\ldots{}
  (Apéry's constant) is the most important test.
\item
  \textbf{New proof of the Basel problem} (conditional). If both α = 1/6
  (Conjecture 2) and R = −π (Conjecture 1) are proved from carry arithmetic,
  then ζ(2) = αR² = π²/6 follows. This depends on two currently unproved
  conjectures; neither is close to resolution.
\end{enumerate}

\begin{center}\rule{0.5\linewidth}{0.5pt}\end{center}

\newpage

\subsection{References}\label{references}

\begin{enumerate}
\def\labelenumi{\arabic{enumi}.}
\tightlist
\item
  P. Diaconis and J. Fulman, \emph{Carries, shuffling, and symmetric functions},
  Adv. Appl. Math. \textbf{43} (2009), 176--196.
\item
  P. Diaconis and J. Fulman, \emph{Carries, shuffling, and an amazing matrix},
  Amer. Math. Monthly \textbf{116} (2009), 788--803.
\item
  J. Fulman, \emph{The carries process revisited}, preprint (2023).
\item
  J. M. Holte, \emph{Carries, combinatorics, and an amazing matrix}, Amer. Math.
  Monthly \textbf{104} (1997), 138--149.
\item
  D. E. Knuth, \emph{The Art of Computer Programming, Volume 2: Seminumerical
  Algorithms}, 3rd ed., Addison-Wesley, 1997. (Chapter 4.3: Multiple-precision
  arithmetic.)
\item
  L. Euler, \emph{De summis serierum reciprocarum}, Comment. Acad. Sci. Petrop.
  \textbf{7} (1740), 123--134.
\item
  R. Remmert, \emph{Classical Topics in Complex Function Theory}, Graduate Texts
  in Mathematics \textbf{172}, Springer, 1998. (Chapter 9: The number π.)
\item
  M. V. Berry and J. P. Keating, \emph{The Riemann zeros and eigenvalue
  asymptotics}, SIAM Rev.~\textbf{41} (1999), 236--266.
\item
  L. F. Richardson and J. A. Gaunt, \emph{The deferred approach to the limit},
  Phil. Trans. R. Soc. A \textbf{226} (1927), 299--361.
\item
  I. S. Gradshteyn and I. M. Ryzhik, \emph{Table of Integrals, Series, and
  Products}, 7th ed., Academic Press, 2007. (Entry 1.421.3.)
\item
  P. G. L. Dirichlet, \emph{Beweis des Satzes, dass jede unbegrenzte
  arithmetische Progression\ldots{}}, Abh. Königl. Preuss. Akad. Wiss. (1837),
  45--81.
\item
  G. Tenenbaum, \emph{Introduction to Analytic and Probabilistic Number Theory},
  3rd ed., Graduate Studies in Mathematics \textbf{163}, AMS, 2015. (Chapter
  II.5: Selberg--Delange method.)
\item
  T. Apostol, \emph{Introduction to Analytic Number Theory}, Springer, 1976.
  (Chapters 6, 12: Dirichlet characters and L-functions.)
\item
  H. Weyl, \emph{Über die asymptotische Verteilung der Eigenwerte}, Nachr. Ges.
  Wiss. Göttingen (1911), 110--117.
\item
  D. Zagier, \emph{Values of zeta functions and their applications}, in
  \emph{First European Congress of Mathematics}, Progress in Math. \textbf{120},
  Birkhäuser, 1994, 497--512.
\item
  E. C. Titchmarsh, \emph{The Theory of the Riemann Zeta-Function}, 2nd
  ed.~(revised by D. R. Heath-Brown), Oxford, 1986.
\item
  {[}P2{]} Companion paper: \emph{The sector ratio in binary multiplication:
  from Markov failure to transcendence}.
\item
  {[}G{]} Companion paper: \emph{The angular uniqueness of base 2 in positional
  multiplication}.
\item
  {[}E{]} Companion paper: \emph{The Trace Anomaly of Binary Multiplication}.
  (Identifies conditionally the mechanism for \(R \to -\pi\) via resolvent
  universality, assuming the LMH.)
\item
  {[}A{]} Companion paper: \emph{Spectral theory of carries in positional
  multiplication}. (Foundation: the \(m\)-bit Equidistribution Lemma extending
  Diaconis--Fulman to the transfer operator.)
\item
  {[}F{]} Companion paper: \emph{Exact covariance structure of binary carry
  chains}. (Carry expectation \(E[c_j] = (j-1)/4\), off-diagonal covariance
  \(\text{Cov}(c_j, g_i h_{j-i}) = 1/8\).)
\item
  {[}B{]} Companion paper: ``Carry Polynomials and the Euler Product: An
  Approximation Framework,'' this series.
\end{enumerate}

\end{document}
